\section{What the objective does not guarantee}
\label{sec:theory}

The experiments concern fitted finite networks. The following population statements ask a narrower
question: which guarantees follow from prediction error and a Gaussian latent marginal alone?
They do not characterize optimization or identify the cause of the empirical failures.

\subsection{Non-collapse does not select the controlled state}
\label{sec:theory-rep}

Write observations as $x_t=R(p_t,b)$, where $p_t$ is controlled state and $b$ is an episode-constant
nuisance recoverable from the image. Assume that $b$ is independent of states and actions in the
training data. At deployment, assume it is independent of the initial state and transition randomness,
and does not directly affect dynamics or reward. Let $\mathrm{SIGReg}_\infty$ be the nonnegative
population characteristic-function discrepancy, which vanishes for $\Nrm(0,I_D)$
(Appendix~\ref{app:proofs}).

\begin{proposition}[Gaussian nuisance optimum; extending \citet{sobal2022slow}]
\label{prop:representation}
If the encoder class contains a map $f(R(p,b))=\varphi(b)$ with
$\varphi(b)\sim\Nrm(0,I_D)$ and the predictor class contains the last-latent copy map, their population
prediction loss and $\mathrm{SIGReg}_\infty$ are both zero. This is a global optimum with covariance
$I_D$ and $I(z_t;p_{1:T})=0$ under the training distribution. A deployed controller observing only
these latents, time, and its own actions has the value of a randomized state-blind policy.
\end{proposition}

The argument is simple: nuisance changes across episodes, satisfying the marginal constraint, but
is constant within each episode, making prediction exact. Its deployment value follows by giving a
state-blind policy the same private nuisance draw and controller randomness. The Gaussianization
assumption requires a rich nuisance and a realizable measurable map; a finite texture does not
literally instantiate it. The finite-batch regularizer also differs from the population functional.
A complementary innovation bound in Appendix~\ref{app:proofs} shows that an encoder preserving all
state distances at a fixed scale must incur prediction error for unpredictable state innovations.

\subsection{Decodability does not set a useful prediction scale}
\label{sec:amplitude}

Even preserving all state information leaves its geometry unconstrained. For any $\varepsilon>0$,
we can place the two values of a relevant bit less than $2\varepsilon$ apart while keeping the
\emph{entire} latent distribution exactly standard Gaussian.

\begin{proposition}[Exact information at arbitrarily small scale]
\label{prop:amplitude}
For every $\varepsilon>0$ and $D\ge1$, there is a measurable observation process with episode-constant
nuisance $b$ and a fresh independent fair state bit $p_t$, and an encoder $f_\varepsilon$, such that
$z_t\sim\Nrm(0,I_D)$ and both $p_t$ and $b$ are exactly recoverable almost surely. For fixed $b$,
$0<\|f_\varepsilon(R(1,b))-f_\varepsilon(R(0,b))\|<2\varepsilon$.
The optimal next-latent predictor is independent of the next bit and has coordinate-mean squared
error in $(0,\varepsilon^2/D]$. With reward $\mathbf1\{a_t=p_t\}$, latent feedback attains value $1$
per step, while every state-blind policy attains $1/2$.
\end{proposition}

The construction partitions one Gaussian coordinate into narrow intervals and uses the bit to
select the lower or upper half of each interval's conditional probability mass. The nuisance fixes
the interval and the within-half quantile. This preserves the Gaussian marginal and exact decoding
while making the two states arbitrarily close. The full proof is in Appendix~\ref{app:proofs}.
It uses infinite-precision measurable coding; the next bit is intrinsically unpredictable. It therefore
separates information from scale, without proving avoidable prediction error or explaining deletion
of the predictable gate wall.

\paragraph{Why scale matters for a downstream read-out.}
For any fixed $L$-Lipschitz downstream score $q$ (for example, an action-value read-out)
and any coupled real and predicted next latents,
\begin{equation}
\label{eq:readout}
\left|\E q(Z')-\E q(\widehat Z')\right|
\le L\sqrt{\E\|Z'-\widehat Z'\|^2}.
\end{equation}
This standard bound follows from Lipschitz continuity and Cauchy--Schwarz. Prediction error controls
score error only relative to the read-out's sensitivity. Exact decodability alone supplies no bound
on that sensitivity. We do not estimate $L$ or claim a control guarantee for the fitted networks;
the inequality makes explicit the missing bridge between a geometric loss and a decision.
